\chapter{The Halo Estimate of Liu--Wan--Xiao}
\label{chap:lwx-halo}

Chapters~\ref{chap:qmf}--\ref{chap:jacobs} computed slopes at one weight at a time.  Weight space
$\mathcal{W}$ is a union of open discs, and the picture above it changes character near the
boundary: Liu, Wan and Xiao \cite{LWX} proved, for the spectral curve of overconvergent
automorphic forms on a definite quaternion algebra, that over the boundary annulus it is a
disjoint union of ``pieces'' finite and flat over the annulus, on each of which the slope of
$a_p$ is a fixed multiple of $v(T)$ --- the multiple being independent of the weight.  This is
the \emph{spectral halo}, an analogue for $\mathrm{Spc}_D$ of the Coleman--Mazur--Buzzard--Kilford
conjecture; the precise statements, and exactly which parts of them are formalised here, are
laid out in Chapter~\ref{chap:lwx-slopes} and \S\ref{sec:what-15-says}.

Their argument has three layers, and this chapter is the first: a purely integral estimate.
Let $\Lhalo = \Zp\llbracket T, pT^{-1}\rrbracket$ be the ring of \cite[Lemma 3.15]{LWX}.  On the
integral model of overconvergent forms the operator $U_p$ has a matrix $P$ with entries in
$\Lhalo$, and \cite[Theorem 3.16]{LWX} says that its Fredholm determinant
$\Char(P) = \sum_n c_n X^n$ exists and satisfies
\[
  c_n \in T^{\lambda(n)}\,\Lhalo, \qquad
  \lambda(n) = \sum_{k<n}\Bigl(\Bigl\lfloor \tfrac{k}{t}\Bigr\rfloor -
  \Bigl\lfloor \tfrac{k}{pt}\Bigr\rfloor\Bigr),
\]
where $t$ is the size of the class set.  Specialising $T$ at a point $T_0$ of the boundary
annulus $p^{-1} < \lVert T_0\rVert < 1$ gives \cite[Corollary 3.18]{LWX}: the Newton polygon of
the specialised series lies on or above an explicit polygon whose slopes are the numbers
$(\lfloor k/t\rfloor - \lfloor k/pt\rfloor)\,v(T_0)$.  The remarkable feature, and the whole
point, is that this lower bound is a \emph{multiple of $v(T_0)$}: it degenerates to the trivial
bound as one moves into the interior, and it is sharp near the boundary.

Everything in this chapter is unconditional and, except where stated, is proved for an arbitrary
combinatorial datum of the shape that \cite[Proposition 3.1]{LWX} produces; the identification of
that datum with the genuine Hecke operator is Chapter~\ref{chap:lwx-model}, and the slope
consequences are Chapter~\ref{chap:lwx-slopes}.  Throughout $p$ is an odd prime, so that
$q = p$ in \cite{LWX}'s notation; $p = 2$ is out of scope.

\section{The integral halo ring}

\begin{definition}
  \label{def:haloint}
  \lean{LWX.HaloInt, LWX.HaloInt.norm_def, LWX.HaloInt.norm_le_one}
  \leanok
  $\Lhalo$ is realised coefficientwise, as \cite[proof of Corollary 3.18]{LWX} describes its
  elements: a two-sided stream $d : \Z \to \Zp$ with $v(d_j) \geq \max(0, -j)$, that is
  $\lVert d_j\rVert \leq p^{\min(0,j)}$, with pointwise addition, convolution product, and the
  gauge norm
  \[
    \lVert d\rVert = \sup_{j\in\Z} \lVert d_j \rVert\, p^{-j} \in [0,1] .
  \]
\end{definition}

Both tails of the convolution die --- the left-hand one by the halo bound, the right-hand one by
integrality --- so the product is a convergent sum over $\Z$ and the ring is complete.

\begin{proposition}
  \label{prop:haloint-norm}
  \uses{def:haloint}
  \lean{LWX.HaloInt.T, LWX.HaloInt.norm_T_mul, LWX.HaloInt.constRingHom,
        LWX.HaloInt.norm_const, LWX.HaloInt.exists_T_pow_mul_of_norm_le,
        LWX.HaloInt.norm_le_of_eq_T_pow_mul}
  \leanok
  The constant embedding $\Zp \to \Lhalo$ is isometric; the element $T = \delta_1$ acts as the
  coefficient shift and satisfies $\lVert Tx\rVert = p^{-1}\lVert x\rVert$ exactly.
  Consequently
  \[
    x \in T^{k}\,\Lhalo \iff \lVert x\rVert \leq p^{-k},
  \]
  which is \cite[Lemma 3.15]{LWX} in the only form the estimate consumes.
\end{proposition}

Divisibility by a power of $T$ has become a norm inequality; this is what makes the halo estimate
an instance of the general Hadamard bound of Chapter~\ref{chap:compact} rather than an argument
about ideals.

\begin{definition}[Specialisation at a halo point]
  \label{def:halo-specialize}
  \uses{def:haloint}
  \lean{LWX.HaloInt.specialize, LWX.HaloInt.specializeHom, LWX.HaloInt.continuous_specialize,
        LWX.HaloInt.norm_specialize_le}
  \leanok
  Let $K$ be a complete ultrametric field, $\psi : \Zp \to K$ an isometric ring homomorphism and
  $T_0 \in K$ with $p^{-1} < \lVert T_0\rVert < 1$.  Then $\sum_{j\in\Z}\psi(d_j)T_0^{\,j}$
  converges, and $d \mapsto d(T_0)$ is a continuous ring homomorphism $\Lhalo \to K$
  satisfying $\lVert d(T_0)\rVert \leq \lVert T_0\rVert^{k}$ whenever $\lVert d\rVert\leq p^{-k}$.
\end{definition}

The two hypotheses on $\lVert T_0\rVert$ are exactly what makes the two tails converge; this is
the annulus $\mathcal{W}^{>1/p}$ of \cite{LWX}.  Continuity is needed because $\Char(P)$ is a
limit of determinants of minors, and the specialisation has to be pushed through that limit
(\S\ref{sec:seam}).

\begin{proposition}[{The Banach--Tate ring $A = \Lhalo[1/T]$}]
  \label{prop:halotate}
  \uses{def:haloint, def:istate}
  \lean{LWX.HaloTate, LWX.HaloTate.ofIntRingHom, LWX.HaloTate.T,
        LWX.HaloTate.norm_T_mul, LWX.HaloTate.pseudoUniformizer, LWX.HaloTate.norm_le_one_iff,
        LWX.HaloTate.exists_T_zpow_mul_ofInt}
  \leanok
  Inverting $T$ gives the ring $A$ of streams satisfying a \emph{shifted} halo bound.  It is
  complete for the gauge norm, contains $\Lhalo$ isometrically as its unit ball, and $T$ is a
  unit with $\lVert Tx\rVert = p^{-1}\lVert x\rVert$: a multiplicative pseudo-uniformizer.  So
  $A$ is a Banach--Tate ring in the sense of Definition~\ref{def:istate}, and every element is
  $T^{-k}$ times an integral one.
\end{proposition}

No Noetherian hypothesis is claimed, and none is available: this is precisely why
Chapter~\ref{chap:riesz-coleman} was developed Noetherian-free.

\section{The universal character}

\begin{definition}
  \label{def:univchar}
  \uses{def:haloint}
  \lean{LWX.teichmuller, LWX.oneUnitPart, LWX.qlog, LWX.qlog_mul, LWX.oneAddTPow,
        LWX.univChar, LWX.univChar_mul, LWX.norm_univChar_le}
  \leanok
  Write $\Zp^{\times} = \Delta \times (1+p\Zp)$, $x = \omega_{\mathrm{T}}(x)\langle x\rangle$,
  with the Teichm\"uller lift $\omega_{\mathrm{T}}(x) = \lim_n x^{p^n}$; let
  $\ell\langle x\rangle = \tfrac1p\log\langle x\rangle \in \Zp$.  For a nebentypus
  $\omega : (\Z/p)^{\times}\to\Zp^{\times}$ the \emph{universal character} of
  \cite[Notation 2.1]{LWX} is
  \[
    [x] \;=\; \omega(\bar x)\,(1+T)^{\ell\langle x\rangle} \;\in\; (\Lhalo)^{\times},
    \qquad (1+T)^{s} = \sum_{r\geq 0}\binom{s}{r}T^{r} .
  \]
  It is multiplicative and takes values in the unit ball.
\end{definition}

Mathlib's \texttt{WittVector.teichmuller} is a different map, so the lift, the logarithm on
$1+p\Zp$ and its additivity are all built here; additivity of $\log$ is what allows
$\log((cz+d)/d_0)/p$ to be split into a constant plus a series, and that split is what makes the
binomial estimates of the next section applicable.

\section{Tilted degree}

\cite{LWX} call \S3.3--\S3.13 ``a mild $p$-adic analysis computation (which is the core of our
paper)''.  Its object is a notion of degree for functions $\Zp\to\Zp$ that is stable under the
operations the matrix entries are built from.

\begin{definition}
  \label{def:tilteddeg}
  \lean{LWX.TiltedDeg}
  \leanok
  A function $f : \Zp\to\Zp$ has \emph{tilted degree $\leq n$} if
  $\lVert \tilde\Delta^{m}f(z)\rVert \leq p^{\,n-m}$ for all $m$ and all $z$, where
  $\tilde\Delta$ is the forward difference $f(z+1)-f(z)$.
\end{definition}

This is clause (1) of \cite[Definition-Proposition 3.8]{LWX}.  Taking the bound
\emph{everywhere} rather than only at $z=0$ costs nothing and buys two things: no Mahler
expansion is needed to state it, and the class is closed under pointwise limits.  Clause (2),
the Mahler-coefficient form, is then never needed and is not stated.

\begin{lemma}[The calculus]
  \label{lem:tilted-calculus}
  \uses{def:tilteddeg}
  \lean{LWX.TiltedDeg.mul, LWX.TiltedDeg.of_tendsto, LWX.TiltedDeg.mono, LWX.TiltedDeg.sum,
        LWX.tiltedDeg_const, LWX.fwdDiff_iter_mul, LWX.fwdDiff_iter_choose}
  \leanok
  Tilted degree is monotone in $n$, additive under products, closed under convergent sums and
  pointwise limits, and constants have tilted degree $0$.  The engine is the Leibniz rule
  $\tilde\Delta^{m}(fg) = \sum_{i+j=m}\tilde\Delta^{i}f \cdot (\tilde\Delta^{j}g)(\cdot + i)$,
  which is \cite[(3.5.1)]{LWX}, together with
  $\tilde\Delta\binom{z}{n} = \binom{z}{n-1}$, which is \cite[(3.5.2)]{LWX}.
\end{lemma}

The Chu--Vandermonde input \cite[(3.5.3)]{LWX} is mathlib's binomial-ring
\texttt{add\_choose\_eq}, and the Mahler coefficient $a_m(f) = \tilde\Delta^{m}f(0)$ is mathlib's
\texttt{fwdDiff\_iter\_eq\_sum\_shift}.

\begin{lemma}[{\cite[Lemmas 3.12 and 3.13]{LWX}}]
  \label{lem:tilted-binomial}
  \uses{def:tilteddeg, lem:tilted-calculus}
  \lean{LWX.tiltedDeg_choose_of_pMul, LWX.tiltedDeg_choose_monomial_p,
        LWX.tiltedDeg_choose_of_logShape}
  \leanok
  If $f : \Zp\to p\Zp$ then $z\mapsto\binom{f(z)}{n}$ has tilted degree $\leq \lfloor n/p\rfloor$;
  and if $g(z) = \sum_k p^{k-1}(a_k/k)z^{k}$ has the shape of a normalised logarithm then
  $z\mapsto\binom{g(z)}{r}$ has tilted degree $\leq 0$.
\end{lemma}

These two estimates are where the odd prime is used: the proof rests on the Legendre margin
$2\,v_p(n!)\leq n-1$ and on $4v\leq p^{v}+3$, both of which fail at $p=2$.

\section{The matrix of $U_p$}

\begin{definition}
  \label{def:localmat}
  \lean{LWX.LocalMat, LWX.LocalMat.IsUpShape}
  \leanok
  A \emph{local matrix} is $\delta = \begin{pmatrix} a & b\\ c& d\end{pmatrix}\in M_2(\Zp)$ with
  $p \mid c$ and $d \in \Zp^{\times}$: the monoid $M_1$ of \cite[(2.3.3)]{LWX}.  It has
  \emph{$U_p$-shape} if in addition $a \in p\Zp$, which is clause (3) of
  \cite[Proposition 3.1]{LWX}.
\end{definition}

\begin{definition}[The entry stream, {\cite[Propositions 3.4 and 3.14]{LWX}}]
  \label{def:entry}
  \uses{def:localmat, def:univchar, def:haloint}
  \lean{LWX.entry, LWX.entryCoeff, LWX.coeff_entry}
  \leanok
  \cite[Proposition 3.4]{LWX} computes the matrix of a single $\Vert_\delta$ in the Mahler basis:
  \[
    P_{m,n}(\delta) \;=\; \tilde\Delta^{m}\Bigl(\;\binom{\tfrac{az+b}{cz+d}}{n}\,[\,cz+d\,]\;
    \Bigr)\Big|_{z=0} .
  \]
  Expanding $[cz+d] = [d_0](1+T)^{g(z)}$ with $g(z) = \tfrac1p\log\bigl((cz+d)/d_0\bigr)$ gives
  the $T$-coefficient stream
  $r \mapsto [d_0]\,\tilde\Delta^{m}\bigl(\binom{f(z)}{n}\binom{g(z)}{r}\bigr)\big|_{z=0}$,
  and we take that stream as the definition of the entry $P_{m,n}(\delta)\in\Lhalo$.
\end{definition}

That the stream computes the action on the integral model is a separate statement, proved in
Chapter~\ref{chap:lwx-model} (Proposition~\ref{prop:prop34}); here it is simply the matrix we
estimate.

\begin{proposition}[{\cite[Proposition 3.14(1)]{LWX}}]
  \label{prop:entry-bound}
  \uses{def:entry, lem:tilted-binomial, prop:haloint-norm}
  \lean{LWX.norm_entry_le, LWX.norm_entryCoeff_le, LWX.norm_entry_le_M1}
  \leanok
  For a local matrix of $U_p$-shape,
  \[
    \bigl\lVert P_{m,n}(\delta)\bigr\rVert \;\leq\; p^{-(m - \lfloor n/p\rfloor)},
    \qquad\text{i.e.}\qquad
    P_{m,n}(\delta) \in \mathfrak{m}^{\max(m-\lfloor n/p\rfloor,\,0)} .
  \]
  Without the shape hypothesis, that is for a general element of $M_1$, one still has
  $\lVert P_{m,n}(\delta)\rVert \leq p^{-(m-n)}$.
\end{proposition}

\begin{proof}
  \leanok
  \uses{lem:tilted-binomial, lem:tilted-calculus, def:univchar}
  The $U_p$-shape makes $f(z) = (az+b)/(cz+d)$ take values in $p\Zp$, so
  $\binom{f(z)}{n}$ has tilted degree $\leq\lfloor n/p\rfloor$ by the first half of
  Lemma~\ref{lem:tilted-binomial}; $g$ is a normalised logarithm, so $\binom{g(z)}{r}$ has tilted
  degree $\leq 0$ by the second half; the product has tilted degree
  $\leq \lfloor n/p\rfloor$ by Lemma~\ref{lem:tilted-calculus}, and evaluating the definition of
  tilted degree at the $m$-th difference and $z=0$ is the claim.  The constant $[d_0]$ is a unit.
\end{proof}

\begin{definition}
  \label{def:updatum}
  \uses{def:localmat, def:entry}
  \lean{LWX.UpDatum, LWX.UpDatum.matrix, LWX.UpDatum.op, LWX.UpDatum.matrixCoeff_op}
  \leanok
  A \emph{$U_p$-datum} on a class set $\iota$ consists of, for each $i\in\iota$, exactly $p$
  summands, each given by a target index $\mathrm{tgt}(i,j)\in\iota$ and a local matrix
  $\mathrm{mat}(i,j)$ of $U_p$-shape.  This is the content of \cite[Proposition 3.1]{LWX} that
  the estimate consumes; the column count of clause (2) is not needed.  Assembling the entry
  streams block by block gives a matrix on $\iota\times\N$ and, by cofinite column decay, a
  continuous operator $D.\mathrm{op}$ on $c(\iota\times\N, \Lhalo)$.
\end{definition}

\section{Theorem 3.16 and Corollary 3.18}

\begin{definition}
  \label{def:lwxlambda}
  \lean{LWX.lwxLambda, LWX.lwxLambda_succ, LWX.lwxLambda_eq_sum_comp}
  \leanok
  With $t = \lvert\iota\rvert$, set
  $\lambda(n) = \sum_{k<n}\bigl(\lfloor k/t\rfloor - \lfloor k/pt\rfloor\bigr)$.
\end{definition}

\begin{theorem}[{\cite[Theorem 3.16]{LWX}}]
  \label{thm:lwx316}
  \uses{def:updatum, def:lwxlambda, lem:two-sided, prop:entry-bound, def:iscompactoid}
  \lean{LWX.norm_charCoeff_upOp_le, LWX.exists_charCoeff_upOp_eq_T_pow_mul,
        LWX.summable_minor_upOp}
  \leanok
  For every $U_p$-datum $D$ on a class set of size $t$ and every nebentypus $\omega$, the
  Fredholm determinant $\Char(P) = \sum_n c_n X^n$ of $D.\mathrm{op}$ is defined over $\Lhalo$
  and
  \[
    \lVert c_n\rVert \;\leq\; p^{-\lambda(n)},
    \qquad\text{equivalently}\qquad
    c_n \in T^{\lambda(n)}\,\Lhalo .
  \]
\end{theorem}

\begin{proof}
  \leanok
  \uses{lem:two-sided, prop:entry-bound, def:lwxlambda, prop:haloint-norm}
  Proposition~\ref{prop:entry-bound} says the matrix entry at $\bigl((i,m),(j,n)\bigr)$ has norm
  at most $p^{-\max(m-\lfloor n/p\rfloor,\,0)}$: a two-sided weight, with row weight
  $w(i,m) = m$ and column weight $w'(j,n) = \lfloor n/p\rfloor$.  Lemma~\ref{lem:two-sided}
  therefore bounds $\lVert c_n\rVert$ by $p^{-f(n)}$ for any $f$ with
  $f(n)\leq \sum_{a\in S}w(a) - \sum_{a\in S}w'(a)$ over $n$-element $S$, and the combinatorial
  content is that $\lambda$ is such an $f$: on $\iota\times\N$ with $\lvert\iota\rvert = t$, the
  smallest possible value of $\sum_{(i,m)\in S}\bigl(m - \lfloor m/p\rfloor\bigr)$ over
  $\lvert S\rvert = n$ is attained by the first $n$ indices in the block order, which gives
  $\sum_{k<n}\bigl(\lfloor k/t\rfloor - \lfloor \lfloor k/t\rfloor/p\rfloor\bigr) = \lambda(n)$.
  Summability of the minors is the same estimate.  Finally $T$-divisibility is
  Proposition~\ref{prop:haloint-norm}.
\end{proof}

\cite{LWX} obtain the estimate by conjugating $P$ by $\mathrm{diag}(1,\dots,T,\dots,T^2,\dots)$
and reading off column divisibility.  Lemma~\ref{lem:two-sided} performs the same computation at
the level of minors, which are invariant under diagonal conjugation, so no $T^{-1}$ ever appears
and the whole argument stays inside the integral ring.

\begin{definition}
  \label{def:speccharseries}
  \uses{thm:lwx316, def:halo-specialize}
  \lean{LWX.specCharSeries, LWX.specCharSeries_coeff_zero, LWX.lwxSlopes,
        LWX.monotone_lwxSlopes}
  \leanok
  At a halo point $T_0$ with $p^{-1}<\lVert T_0\rVert<1$ set
  $\Char(P)(T_0) = \sum_n c_n(T_0)X^n \in K\llbracket X\rrbracket$.  Its constant coefficient is
  $1$.  The comparison slope sequence is
  $s_k = \bigl(\lfloor k/t\rfloor - \lfloor k/pt\rfloor\bigr)\cdot v(T_0)$, which is monotone.
\end{definition}

\begin{theorem}[{\cite[Corollary 3.18]{LWX}}]
  \label{thm:lwx318}
  \uses{def:speccharseries, thm:lwx316, def:ofslopes, def:isnewtonpolygonof, thm:ofslopes}
  \lean{LWX.norm_specCharSeries_coeff_le, LWX.isBelow_newtonPolygon_specCharSeries}
  \leanok
  For every $T_0$ in the boundary annulus, $v\bigl(c_n(T_0)\bigr) \geq \lambda(n)\,v(T_0)$, and
  the Newton polygon of $\Char(P)(T_0)$ lies on or above the polygon anchored at the origin
  whose $k$-th unit slope is $\bigl(\lfloor k/t\rfloor - \lfloor k/pt\rfloor\bigr)v(T_0)$.
\end{theorem}

\begin{proof}
  \leanok
  \uses{thm:lwx316, def:halo-specialize, thm:ofslopes, def:isnewtonpolygonof}
  The coefficient bound is Theorem~\ref{thm:lwx316} pushed through the specialisation estimate of
  Definition~\ref{def:halo-specialize}.  For the polygon statement, the constant coefficient is
  $1$ so the polygon is anchored at the origin, all coefficients have norm at most $1$ so the
  valuation sequence is admissible, and $\sum_{i<k}s_i = \lambda(k)\,v(T_0) \leq
  v(c_k(T_0))$ is the coefficient bound after taking logarithms; now apply the maximality clause
  of Definition~\ref{def:isnewtonpolygonof} to the comparison polygon of
  Definition~\ref{def:ofslopes}.
\end{proof}

\section{Sharpness}

The lower bound of Theorem~\ref{thm:lwx318} is attained exactly when a single integral
coefficient is a unit.  Write $c_n = \sum_m b_{n,m}T^{m}$ for the $T$-expansion of the integral
coefficient.

\begin{theorem}[The equality clause of {\cite[Corollary 3.18]{LWX}}]
  \label{thm:lwx-sharp}
  \uses{thm:lwx318, def:speccharseries}
  \lean{LWX.norm_specCharSeries_coeff_eq_iff, LWX.norm_specCharSeries_coeff_eq_of_isUnit,
        LWX.norm_specCharSeries_coeff_le_of_not_isUnit}
  \leanok
  $v\bigl(c_n(T_0)\bigr) = \lambda(n)v(T_0)$ if and only if $b_{n,\lambda(n)} \in \Zp^{\times}$.
  Otherwise there is a definite margin:
  \[
    v\bigl(c_n(T_0)\bigr) \;\geq\; \lambda(n)\,v(T_0) + \min\bigl(v(T_0),\, 1-v(T_0)\bigr).
  \]
\end{theorem}

\begin{proof}
  \leanok
  \uses{thm:lwx318}
  The term-by-term estimate \cite[(3.18.1)]{LWX} shows that in
  $c_n(T_0) = \sum_m \psi(b_{n,m})T_0^{\,m}$ the term $m = \lambda(n)$ strictly dominates all
  others when $b_{n,\lambda(n)}$ is a unit: terms with $m < \lambda(n)$ are smaller because
  integrality forces $v(b_{n,m}) \geq \lambda(n)-m$, and terms with $m > \lambda(n)$ are smaller
  because $\lVert T_0\rVert<1$.  An ultrametric sum with a unique dominant term has that term's
  norm.  If $b_{n,\lambda(n)}$ is not a unit then every term loses either a factor
  $\lVert T_0\rVert$ or a factor $p^{-1}\lVert T_0\rVert^{-1}$, which is the stated margin.
\end{proof}

The if-and-only-if is the hinge of Chapter~\ref{chap:lwx-slopes}: it converts a geometric
statement about the Newton polygon at one halo point into the $T$-free arithmetic statement that
a particular integral coefficient is a unit --- and a $T$-free statement can then be transported
to \emph{every} halo point.

\section{Riesz theory at a halo vertex}
\label{sec:tate-riesz}

Over the Banach--Tate ring $A$ of Proposition~\ref{prop:halotate} the halo $U_p$ is an honest
compact operator, and Chapter~\ref{chap:riesz-coleman} applies to it.

\begin{definition}
  \label{def:tateop}
  \uses{prop:halotate, def:updatum, def:iscompactoid}
  \lean{LWX.UpDatum.tateMatrix, LWX.UpDatum.tateOp, LWX.UpDatum.isCompactoid_tateOp}
  \leanok
  Conjugate the matrix $P$ by $\mathrm{diag}(T^{\lfloor n/t\rfloor - \lfloor n/pt\rfloor})$, as
  \cite[proof of Theorem 3.16]{LWX} does, and transpose.  The result is a compactoid operator on
  $c(\iota\times\N, A)$.
\end{definition}

The transpose is not cosmetic: after rescaling it is the \emph{columns} of $P'$ that decay
uniformly, so the honest compact operator in the sense of Definition~\ref{def:iscompactoid} ---
which asks for row decay --- is $(P')^{\mathsf{T}}$.

\begin{proposition}
  \label{prop:tateop-det}
  \uses{def:tateop, lem:diag-intertwine}
  \lean{LWX.UpDatum.minor_tateOp, LWX.UpDatum.charPowerSeries_tateOp,
        LWX.UpDatum.charCoeff_tateOp}
  \leanok
  Principal minors are invariant under transposition and diagonal conjugation, so the Fredholm
  determinant of the operator of Definition~\ref{def:tateop} is $\Char(P)$ read in
  $A\llbracket X\rrbracket$.
\end{proposition}

\begin{definition}
  \label{def:halovertex}
  \uses{def:tateop, thm:lwx-sharp}
  \lean{LWX.UpDatum.IsHaloVertex, LWX.HaloInt.isUnit_of_isUnit_coeff_zero,
        LWX.UpDatum.isHaloVertex_of_isUnit_coeff}
  \leanok
  An index $n$ is a \emph{halo vertex} if $c_n = T^{\lambda(n)}\cdot(\text{unit})$ in $\Lhalo$.
  This is the situation of \cite[Remark 3.25]{LWX}; by Theorem~\ref{thm:lwx-sharp} it is implied
  by $b_{n,\lambda(n)}$ being a unit, that is, by the Newton polygon touching the lower bound at
  $n$.
\end{definition}

\begin{theorem}[Riesz decomposition at a halo vertex]
  \label{thm:halo-riesz}
  \uses{def:halovertex, thm:slope-factorisation, thm:riesz-coleman, prop:tateop-det}
  \lean{LWX.UpDatum.isDominantIndex_charPowerSeries_tateOp,
        LWX.UpDatum.isMultiplicative_charCoeff_tateOp,
        LWX.UpDatum.exists_isDominantFactorization_tateOp,
        LWX.UpDatum.exists_rieszColemanProjection_tateOp}
  \leanok
  At a halo vertex $n$, the index $n$ is dominant for $\Char(P)$ at radius
  $p^{\lambda(n)-\lambda(n-1)}$ and the dominant coefficient is a multiplicative unit.  Hence
  $\Char(P) = P_n\,G$ with $\deg P_n = n$ (Theorem~\ref{thm:slope-factorisation}), and
  Theorem~\ref{thm:riesz-coleman} gives the corresponding decomposition of
  $c(\iota\times\N, A)$: a $U_p$-stable summand of rank $n$ carrying exactly the first $n$
  slopes, over the whole boundary annulus at once.
\end{theorem}

\begin{proof}
  \leanok
  \uses{def:halovertex, thm:lwx316, thm:slope-factorisation, thm:riesz-coleman}
  Dominance at the stated radius is convexity of $\lambda$: the increments
  $\lfloor k/t\rfloor - \lfloor k/pt\rfloor$ are non-decreasing, so $\lambda(k) - k\nu$ is
  minimised at $k = n$ for $\nu$ the $n$-th increment, and the inequality is strict beyond $n$
  because the increments grow.  The coefficient bound of Theorem~\ref{thm:lwx316} gives the
  inequality for $k \neq n$ and the vertex hypothesis gives equality at $n$ together with
  multiplicativity of the norm on $c_n$, since $\lVert T\rVert$ is multiplicative and units of
  $\Lhalo$ are detected by their leading coefficient.  Now apply
  Theorem~\ref{thm:slope-factorisation} and then Theorem~\ref{thm:riesz-coleman}.
\end{proof}

\cite[Remark 3.25]{LWX} appeals to \cite[Proposition 3.2.2]{Kedlaya} at exactly this point; the
hypothesis that $n$ is a halo vertex is derived in \cite{LWX} from the touching argument of
Theorem 1.3, which is \S\ref{sec:stepone} below, so we carry it as an explicit hypothesis.
